% ======================================================================
% Chapter 20: Gauge Theories with Spontaneous Symmetry Breaking
% (book pages 689-728; PDF pages 712-753)
% ======================================================================
\chapter{自发对称性破缺的规范理论}

在本书的进程中，我们讨论了对称性在量子场论中得以实现的三种不同方式。最简单的情形是显式的全局对称性，它导致相互作用受限制的粒子多重态。第二种可能是自发破缺的全局对称性。于是，如第~11 章所讨论的，%
\footnote{第~11.1 节是本章必要的背景知识，但第~11 章的其余部分则不是。}对称性流仍然守恒，相互作用同样受到限制，但真空态并不尊重该对称性，粒子也不形成明显的对称性多重态。相反，这样的理论包含无质量粒子——戈德斯通玻色子，自发破缺对称性的每一个生成元对应一个。第三种情形是局域对称性，或称规范对称性。如我们在第~15 章所见，这样的对称性要求每一个对称性生成元都存在一个无质量矢量场，并且这些场之间的相互作用受到高度限制。

现在我们很自然地考虑第四种可能性：如果在同一个理论中同时包含局域规范不变性和自发对称性破缺，会发生什么？在本章和下一章中，我们将发现，这些要素的组合为量子场论模型的构造带来新的可能性。我们将看到，自发对称性破缺要求规范矢量玻色子获得质量。然而，这些有质量玻色子的相互作用仍受底层规范对称性的约束，而这些约束可以产生可观测的后果。

在基本粒子物理中，自发破缺的局域对称性的主要应用在于目前被接受的弱相互作用模型。这一模型由 Glashow、Weinberg 和 Salam 提出，将在第~20.2 节介绍。在那里我们将看到，它对弱相互作用现象作出了一系列精确且成功的预言。值得注意的是，这一模型还把弱相互作用与电磁相互作用统一在了一个更大的规范理论之中。

\section{希格斯机制}

在本节中，我们分析一些具有自发对称性破缺的规范理论的简单例子。我们从阿贝尔规范理论开始，然后研究几个非阿贝尔模型的例子。

\subsection{一个阿贝尔例子}

作为我们的第一个例子，考虑一个既与自身又与电磁场耦合的复标量场：
\begin{equation}
  \mathcal{L} = -\frac{1}{4}(F_{\mu\nu})^2 + |D_\mu\phi|^2 - V(\phi),
  \label{eq:20.1}
\end{equation}
其中 $D_\mu = \partial_\mu + ieA_\mu$。该拉格朗日量在局域 $U(1)$ 变换下不变
\begin{equation}
  \phi(x) \to e^{i\alpha(x)}\phi(x),\qquad
  A_\mu(x) \to A_\mu(x) - \frac{1}{e}\partial_\mu\alpha(x) .
  \label{eq:20.2}
\end{equation}
如果我们把 $\mathcal{L}$ 中的势选为如下形式
\begin{equation}
  V(\phi) = -\mu^2\phi^{*}\phi + \frac{\lambda}{4}(\phi^{*}\phi)^2 ,
  \label{eq:20.3}
\end{equation}
其中 $\mu^2 > 0$，场 $\phi$ 将获得真空期望值，$U(1)$ 全局对称性将被自发破缺。该势的极小值出现在
\begin{equation}
  \langle\phi\rangle = \phi_0 = \Bigl( \frac{\mu^2}{\lambda} \Bigr)^{1/2} ,
  \label{eq:20.4}
\end{equation}
或在由 $U(1)$ 对称性 \eqref{eq:20.2} 联系起来的任何其他值处。

让我们把拉格朗日量 \eqref{eq:20.1} 在真空态 \eqref{eq:20.4} 附近展开。把复场 $\phi(x)$ 分解为
\begin{equation}
  \phi(x) = \phi_0 + \frac{1}{\sqrt{2}}\bigl( \phi_1(x) + i\phi_2(x) \bigr) .
  \label{eq:20.5}
\end{equation}
势 \eqref{eq:20.3} 可改写为
\begin{equation}
  V(\phi) = \frac{1}{2}(2\mu^2)\phi_1^2 + \lambda\phi_0\,\phi_1^3
    + \mathcal{O}(\phi^4) ,
  \label{eq:20.6}
\end{equation}
于是场 $\phi_1$ 获得质量 $m = \sqrt{2}\mu$，而 $\phi_2$ 是无质量戈德斯通玻色子。到目前为止，整个分析都遵循第~11.1 节。但现在考虑 $\phi$ 的动能项如何变换。代入展开式 \eqref{eq:20.5}，我们改写
\begin{equation}
  |D_\mu\phi|^2 = \frac{1}{2}(\partial_\mu\phi_1)^2 + \frac{1}{2}
    (\partial_\mu\phi_2)^2 + \sqrt{2}e\phi_0\cdot A^\mu(\partial_\mu\phi_2)
      + \cdots ,
  \label{eq:20.7}
\end{equation}
其中我们略去了在 $A^\mu$、$\phi_1$ 和 $\phi_2$ 中三次和四次的项。\eqref{eq:20.7} 中显式写出的最后一项是光子质量项
\begin{equation}
  \Delta\mathcal{L} = \frac{1}{2}m_A^2\,A_\mu A^\mu ,
  \label{eq:20.8}
\end{equation}
其中质量
\begin{equation}
  m_A = \sqrt{2}e\phi_0
  \label{eq:20.9}
\end{equation}
来自 $\phi$ 的非零真空期望值。注意这一质量项的符号是正确的；$A^\mu$ 的物理类空分量在 \eqref{eq:20.8} 中以下列形式出现
\begin{equation}
  \Delta\mathcal{L} = -\frac{1}{2}m_A^2\,\mathbf{A}^2 ,
\end{equation}
其符号对势能项来说是正确的。

在第~7 章，以及在非阿贝尔情形下的第~16 章，我们论证过，规范玻色子不能获得质量，除非这一质量项与真空极化振幅中的极点相联系。在二维时空中存在这一结果的反例；在那里，如我们在第~19.1 节所见，所需形式的极点可以来自无质量费米子对产生的红外奇点。然而，在四维中，真空极化振幅中的极点只能由无质量标量粒子产生。通常，在对称性未破缺的情形下，不存在这样的粒子。

然而，具有自发破缺的连续对称性的模型必须有无质量戈德斯通玻色子。这些标量粒子具有对称性流的量子数，因而恰好具有作为真空极化中间态出现所需的量子数。在我们现在讨论的模型中，我们可以通过以下方式显式地看到这一极点出现：式~\eqref{eq:20.7} 中的第三项把规范玻色子直接耦合到戈德斯通玻色子 $\phi_2$；这给出如下形式的顶角
\begin{equation}
  \text{($\phi_2$-规范场顶角)} = i\sqrt{2}e\phi_0(-ik^\mu)
  = im_A\,k^\mu .
  \label{eq:20.10}
\end{equation}
如果我们还把质量项 \eqref{eq:20.8} 作为微扰论中的顶角来处理，那么对真空极化振幅的最低阶贡献给出表达式
\begin{equation}
  \Pi^{\mu\nu}(k) = im_A^2\,g^{\mu\nu} + \bigl( m_A k^\mu \bigr)\,
    \frac{i}{k^2}\,\bigl( -m_A k^\nu \bigr)
  = im_A^2\Bigl( g^{\mu\nu} - \frac{k^\mu k^\nu}{k^2} \Bigr) .
  \label{eq:20.11}
\end{equation}
戈德斯通玻色子恰好提供所需的极点，使真空极化振幅具有正确的横向性。

尽管戈德斯通玻色子在这一理论中扮演重要的形式角色，但它并不作为独立的物理粒子出现。看清这一点最容易的方法是做出一种特定的规范选择，称为\emph{幺正规范}。利用局域 $U(1)$ 规范对称性 \eqref{eq:20.2}，我们可以选择 $\alpha(x)$，使得 $\phi(x)$ 在每一点 $x$ 处都成为实值。用这一选择，场 $\phi_2$ 就从理论中被消去。拉格朗日量 \eqref{eq:20.1} 变为
\begin{equation}
  \mathcal{L} = -\frac{1}{4}(F_{\mu\nu})^2 + (\partial_\mu\phi_1)^2
    + e^2\phi_1^2A_\mu A^\mu - V(\phi) .
  \label{eq:20.12}
\end{equation}
如果势 $V(\phi)$ 倾向于 $\phi$ 的非零真空期望值，规范场就获得质量；它还保留与剩余物理场 $\phi_1$ 的耦合。

自发对称性破缺为规范玻色子产生质量的这一机制，由 Higgs、Kibble、Guralnik、Hagen、Brout 和 Englert 加以探索并推广到非阿贝尔情形，如今称为\emph{希格斯机制}。然而，这一机制更早就被应用于超导理论。在第~8 章中，我们构造了二级相变的 Landau 描述。为描述超导体，Landau 和 Ginzburg 将该理论耦合到外电磁场；他们得到的正是拉格朗日量 \eqref{eq:20.1}。由于规范场获得非零质量，外电磁场只能穿透超导体到深度 $m_A^{-1}$。这就解释了 Meissner 效应，即所观测到的宏观磁场被超导体排出的现象。

戈德斯通玻色子在希格斯机制中的作用错综复杂，在目前的讨论层面似乎令人费解。我们首先看到，戈德斯通玻色子的参与原则上是规范玻色子获得质量所必需的。接着我们看到，戈德斯通玻色子可以从理论中被形式地消去。然而，我们可能会争辩说，戈德斯通玻色子并没有完全消失。无质量矢量玻色子只有两个物理极化态；我们在第~16 章看到，纵向极化态无法被产生，它出现在形式体系中仅仅是为了抵消其他非物理的贡献。然而，有质量矢量玻色子必须有三个物理极化态：在其静止系中，它是自旋 1 的客体，无法区分横向极化与纵向极化。人们忍不住要说，规范玻色子是通过\emph{吃掉}戈德斯通玻色子来获得其额外自由度的。在第~21.1 节和第~21.2 节中，我们将通过研究自发破缺的规范理论的量子化与规范固定来澄清这一图景。

\subsection{希格斯机制的系统分析}

希格斯机制可以直接推广到具有非阿贝尔规范对称性的系统。推导出标量场真空期望值的一组值导致规范玻色子质量出现的一般关系并不困难。让我们推导出这一关系，然后把它应用于若干例子。

考虑一个标量场系统 $\phi_i$，它出现在对称群 $G$ 下不变的拉格朗日量中，由变换
\begin{equation}
  \phi_i \to \bigl( 1 + i\alpha^a t^a \bigr)_{ij}\phi_j .
  \label{eq:20.13}
\end{equation}
把 $\phi_i$ 写成实值场是方便的，例如把 $n$ 个复场写成 $2n$ 个实场。于是群表示矩阵 $t^a$ 必须是纯虚的，而且由于它们是厄米的，因而必为反对称的。让我们写出
\begin{equation}
  t^a_{ij} = iT^a_{ij} ,
  \label{eq:20.14}
\end{equation}
于是 $T^a$ 是实的且反对称的。

如果我们把对称群 $G$ 提升为局域规范对称性，作用在 $\phi_i$ 上的协变导数为
\begin{equation}
  D_\mu\phi = \bigl( \partial_\mu - igA^a_\mu t^a \bigr)\phi
  = \bigl( \partial_\mu + gA^a_\mu T^a \bigr)\phi .
  \label{eq:20.15}
\end{equation}
于是 $\phi_i$ 的动能项为
\begin{equation}
  |D_\mu\phi|^2 = \frac{1}{2}(\partial_\mu\phi_i)^2
    + gA^{a\mu}\,\partial_\mu\phi_i\,(T^a\phi)_i
    + \frac{1}{2}g^2A^a_\mu A^{b\mu}\,(T^a\phi)_i(T^b\phi)_i .
\end{equation}
现在让 $\phi_i$ 获得真空期望值
\begin{equation}
  \langle\phi\rangle = \phi_0 ,
  \label{eq:20.16}
\end{equation}
并把这些值附近的 $\phi_i$ 展开。式~\eqref{eq:20.15} 中的最后一项含有具有规范玻色子质量结构的项，
\begin{equation}
  \Delta\mathcal{L} = \frac{1}{2}m^2_{ab}A^a_\mu A^{b\mu} ,
  \label{eq:20.17}
\end{equation}
其中质量矩阵为
\begin{equation}
  m^2_{ab} = g^2\,(T^a\phi_0)_i\,(T^b\phi_0)_i .
  \label{eq:20.18}
\end{equation}
该矩阵是半正定的，因为任何基中任何对角元素都具有如下形式
\begin{equation}
  m^2_{aa} = g^2|T^a\phi_0|^2 \ge 0 \qquad \text{(不求和)} .
\end{equation}
因此，一般情况下，所有规范玻色子都将获得正质量。然而，$G$ 的某个特定生成元 $T^a$ 可能使真空保持不变：
\begin{equation}
  T^a\phi_0 = 0 .
\end{equation}
在这种情况下，生成元 $T^a$ 对 \eqref{eq:20.18} 没有贡献，相应的规范玻色子将保持无质量。

与阿贝尔情形一样，规范玻色子传播子从戈德斯通玻色子得到贡献，这是使真空极化振幅成为横向所必需的。为计算这一贡献，我们需要混合规范玻色子与戈德斯通玻色子的顶角。它来自拉格朗日量 \eqref{eq:20.15} 的第二项。当我们代入标量场的真空期望值 \eqref{eq:20.16} 时，这一项变为
\begin{equation}
  \Delta\mathcal{L} = gA^{a\mu}\,\partial_\mu\phi_i\,(T^a\phi_0)_i .
  \label{eq:20.19}
\end{equation}
这一相互作用项并不涉及 $\phi$ 的全部分量——只涉及那些为戈德斯通玻色子的分量。

作为非阿贝尔希格斯机制的一个例子，考虑具有基本（双重态）表示中标量场的 $SU(2)$ 规范理论。真空期望值把 $SU(2)$ 完全破缺，所有三个规范玻色子都获得质量。作为第二个例子，考虑具有伴随（三重态）表示中标量场的 $SU(2)$。在这种情况下，真空期望值留下一个 $U(1)$ 对称性未破缺，因此一个规范玻色子保持无质量，而另外两个获得质量。正如我们的一般分析所示，对应于未破缺对称性生成元的规范玻色子保持无质量。

有趣的是，这一模型既含有质量规范玻色子也含无质量规范玻色子，这些玻色子之间的区别由自发对称性破缺造成。如果我们把有质量玻色子解释为 $W$ 玻色子，把无质量规范玻色子解释为光子，那么人们忍不住要把这一理论解释为弱相互作用与电磁相互作用的统一模型。Georgi 和 Glashow 曾把这一模型作为弱相互作用理论的严肃候选者提出。%
\footnote{H.~Georgi 和 S.~L.~Glashow，\emph{Phys.\ Rev.\ Lett.} \textbf{28}，1494 (1972)。} 然而，自然界选择了另一个不同的模型，我们将在下一节讨论它。

接下来我们转向一个更复杂的例子。考虑具有伴随表示中标量场的 $SU(3)$ 规范理论。$\phi$ 的协变导数具有如下形式
\begin{equation}
  D_\mu\phi^a = \partial_\mu\phi^a + gf^{abc}A^b_\mu\phi^c ,
  \label{eq:20.32}
\end{equation}
于是规范场质量来自这一项
\begin{equation}
  \Delta\mathcal{L} = \frac{1}{2}g^2\bigl( f^{abc}A^b_\mu\phi^c \bigr)^2 .
  \label{eq:20.33}
\end{equation}
我们可以通过定义如下量把它写得更清楚
\begin{equation}
  \Phi = \phi^a t^a ,
  \label{eq:20.34}
\end{equation}
其中 $t^a$ 是表示 $SU(3)$ 生成元的 $3\times3$ 无迹厄米矩阵。利用这一记号与结构常数的定义 \eqref{eq:15.68}，我们可以把质量项 \eqref{eq:20.33} 改写为
\begin{equation}
  \Delta\mathcal{L} = -g^2\,\mathrm{tr}\Bigl[ [t^a,\Phi][t^b,\Phi] \Bigr]\,
    A^a_\mu A^{b\mu} .
  \label{eq:20.35}
\end{equation}
现在让 $\Phi$ 获得真空期望值
\begin{equation}
  \langle\Phi\rangle = \Phi_0 .
  \label{eq:20.36}
\end{equation}
由于 $\Phi_0$ 是无迹厄米矩阵，我们应当通过把它对角化来分析其效应。原则上，$\Phi_0$ 可以有任意三个总和为零的特征值。然而，当人们把显式的势能函数最小化时，常常会发现保持部分原有对称性的期望值。我们将考虑两个例子。

首先，$\Phi_0$ 可能具有取向
\begin{equation}
  \Phi_0 = |\phi|\cdot\frac{1}{\sqrt{15}}\,
    \mathrm{diag}(1, 1, -2) .
  \label{eq:20.37}
\end{equation}
这一矩阵与四个 $SU(3)$ 生成元 $t^1$、$t^2$、$t^3$ 和 $t^8$ 对易。于是，期望值 \eqref{eq:20.37} 把 $SU(3)$ 自发破缺为 $SU(2)\times U(1)$，并使对应于这四个生成元的规范玻色子保持无质量。$SU(3)$ 的其余四个生成元获得质量
\begin{equation}
  m^2 = \bigl( 3g|\phi| \bigr)^2 ,
  \label{eq:20.40}
\end{equation}
这可以通过把这些矩阵代入式~\eqref{eq:20.35} 来检验。

$\Phi_0$ 的另一种可能取向是
\begin{equation}
  \Phi_0 = |\phi|\cdot\frac{1}{\sqrt{6}}\,
    \mathrm{diag}(1, -1, 0) .
  \label{eq:20.41}
\end{equation}
在这种情况下，只有 $t^3$ 和 $t^8$ 与 $\Phi_0$ 对易，因此原来的 $SU(3)$ 对称性被破缺为 $U(1)\times U(1)$。通过代入 \eqref{eq:20.35}，可以确定对应于 $SU(3)$ 其余生成元的规范玻色子获得质量
\begin{equation}
  t^1, t^2:\ m^2 = \bigl( 2g|\phi| \bigr)^2;\qquad
  t^4, t^5, t^6, t^7:\ m^2 = \bigl( g|\phi| \bigr)^2 .
  \label{eq:20.42}
\end{equation}
更大的对称群提供更多样化的对称性破缺模式，以及更复杂的质量矩阵。我们在习题~20.1 中再考虑一个例子。

\subsection{希格斯机制的形式描述}

到目前为止，我们对希格斯机制的研究基于对标量场拉格朗日量与规范场耦合的显式分析。标量场理论提供了具有自发对称性破缺的系统的最简单例子，它们所允许的显式计算有助于直观理解。但对称性也可以通过其他方式破缺。例如，在超导理论中，电磁作用的阿贝尔规范不变性被金属基态中凝聚的电子对破缺。在第~19.3 节中，我们论证过，在夸克质量非常小的近似下，QCD 具有被夸克-反夸克对凝聚自发破缺的全局对称性。在这些例子中，自发对称性破缺是超出微扰论的强相互作用的结果。我们想要理解，这些更一般的自发对称性破缺机制是否也能赋予矢量玻色子质量，如果能，质量又如何计算。

为进行这一分析，我们将需要从前面的讨论中抽象出几个概念。首先，我们将一般地讨论规范玻色子、戈德斯通玻色子与全局对称性流之间的关系。然后，我们将利用这些信息来构造规范玻色子质量矩阵，而不直接使用拉格朗日量。

首先考虑一个具有全局对称性 $G$ 的任意量子场论 $\mathcal{L}_0$。在第~9.6 节中，我们通过对拉格朗日量施加以无穷小参数 $\alpha^a(x)$ 的局域规范变换，推导了对应于 $G$ 对称性的 Noether 流。用常数 $\alpha^a$ 变换应当使 $\mathcal{L}_0$ 保持不变。于是 $\mathcal{L}_0$ 的更一般变分必然具有如下形式
\begin{equation}
  \delta\mathcal{L}_0 = \partial_\mu\alpha^a\,J^{\mu a} ,
  \label{eq:20.43}
\end{equation}
其中 $J^{\mu a}$ 是由 $\mathcal{L}_0$ 的场构成的一组矢量算符。于是变分原理告诉我们
\begin{equation}
  \partial_\mu J^{\mu a} = 0 .
  \label{eq:20.44}
\end{equation}
我们可以把 $J^{\mu a}$ 识别为全局规范对称性的 Noether 流。

我们现在可以把这一全局对称理论耦合到非阿贝尔规范场，把全局对称性提升为局域对称性。到 $g$ 的一阶，新拉格朗日量应当具有如下形式
\begin{equation}
  \mathcal{L} = \mathcal{L}_0 - gA^a_\mu J^{\mu a} + \mathcal{O}(A^2) .
  \label{eq:20.45}
\end{equation}
为检验这一点，注意变换 \eqref{eq:20.43} 在 $g$ 的最低阶抵消了由 $A^a_\mu$ 的规范变换（式~\eqref{eq:15.46}）引起的变分。$A^2$ 阶及更高阶的项一般可以通过安排来抵消规范变换中的高阶项。因此，只含单个规范场插入的矩阵元可以利用原始全局对称理论的 Noether 流的性质来求值。特别要注意，这些流的守恒律（式~\eqref{eq:20.44}）保证了这些矩阵元的 Ward 恒等式得到满足。

如果理论 $\mathcal{L}_0$ 的全局对称性被自发破缺，这一理论将含有戈德斯通玻色子，它们与 Noether 流有特殊的关系。在长波长下，戈德斯通玻色子成为真空的无穷小对称性转动，$Q^a|0\rangle$，其中 $Q^a$ 是与 $J^{\mu a}$ 相关联的全局荷。于是，算符 $J^{\mu a}$ 具有从真空中产生戈德斯通玻色子所需的正确量子数。设 $|\pi^k(p)\rangle$ 表示戈德斯通玻色子态。一般地，会有一个能够产生或湮灭该玻色子的流 $J^{\mu a}$；我们可以把相应的矩阵元参数化为
\begin{equation}
  \langle 0|J^{\mu a}(x)|\pi^k(p)\rangle = -ip^\mu F^{ak}\ e^{-ip\cdot x} ,
  \label{eq:20.46}
\end{equation}
其中 $p^\mu$ 是玻色子在壳动量，$F^{ak}$ 是常数矩阵。当 $a$ 表示未破缺对称性的生成元时，矩阵元 $F^{ak}$ 为零。于是 $F^{ak}$ 的非零矩阵元把自发破缺对称性的流连接到它们相应的戈德斯通玻色子。由于流 $J^{\mu a}$ 是守恒的，我们得到
\begin{equation}
  0 = \partial_\mu\langle 0|J^{\mu a}(x)|\pi^k(p)\rangle
  = -p^2 F^{ak}\,e^{-ip\cdot x} ,
  \label{eq:20.47}
\end{equation}
这意味着具有非零矩阵元 \eqref{eq:20.46} 的玻色子在壳满足 $p^2 = 0$，因而是无质量的。这是戈德斯通定理的另一个证明。%
\footnote{这个论证的一个特例出现在式~\eqref{eq:19.88} 的讨论中。}

由于我们在本节前面考察的标量场理论应当是这一分析的一个特例，我们应当在那里找到式~\eqref{eq:20.46} 所给关系的一个例子。比较式~\eqref{eq:20.15} 和 \eqref{eq:20.45}，我们看到，对标量场理论而言，
\begin{equation}
  J^{\mu a} = \partial^\mu\phi_i\,(T^a\phi)_i ,
  \label{eq:20.48}
\end{equation}
它正是对应于全局规范对称性的 Noether 流。代入真空期望值 \eqref{eq:20.16}，我们得到
\begin{equation}
  J^{\mu a} = \partial^\mu\phi_i\,(T^a\phi_0)_i ,
  \label{eq:20.49}
\end{equation}
这导致如下矩阵元组
\begin{equation}
  \langle 0|J^{\mu a}(x)|\pi^i(p)\rangle = -ip^\mu\,(T^a\phi_0)_i\,
    e^{-ip\cdot x} .
  \label{eq:20.50}
\end{equation}
利用这一关系，我们可以识别出
\begin{equation}
  F^{ai} = (T^a\phi_0)_i
  \label{eq:20.51}
\end{equation}
这是弱耦合标量场理论中希格斯机制的情况。更精确地说，指标 $i$ 遍及标量场 $\phi$ 的全部分量。然而，我们在式~\eqref{eq:20.19} 下方的讨论中看到，\eqref{eq:20.51} 只对为戈德斯通玻色子的分量 $\phi_i$ 非零，并且只对自发破缺的对称性生成元 $a$ 非零。因此，\eqref{eq:20.51} 的非零分量恰好构成结构 \eqref{eq:20.46}。

作为对象 $T^a\phi_0$ 把自发破缺的生成元与戈德斯通玻色子联系起来的方式的一个具体例证，考虑由矢量表示中标量场破缺 $SU(2)$ 对称性的情形。根据几何图像，绕 1 轴的转动把 $\phi$ 的真空期望值倾入 2 方向，绕 2 轴的转动把这个期望值倾入 1 方向，而绕 3 轴的转动使 $\langle\phi\rangle$ 保持不变。于是规范生成元 $T^1$ 和 $T^2$ 被自发破缺，标量场分量 $\phi_2$ 和 $\phi_1$ 是相应的戈德斯通玻色子。这与显式计算 $(T^a\phi_0)_i$ 诸元素的结果相符：利用 $(T^a)_{bc} = \epsilon^{bac}$，我们得到
\begin{equation}
  (T^a\phi_0)_b = \epsilon^{bac}(\phi_0)_c = v\,\epsilon^{ba3} .
  \label{eq:20.52}
\end{equation}
把这一结果代入公式 \eqref{eq:20.50}，我们看到，每个自发破缺对称性的流都产生并湮灭它自己的戈德斯通玻色子。

现在我们可以在这一一般背景下利用这一形式体系来研究希格斯机制的运作。考虑原始理论 $\mathcal{L}_0$ 与 $G$ 的规范玻色子耦合。为看清希格斯机制如何运作，我们必须计算真空极化振幅。Ward 恒等式要求该振幅是横向的，因此它必然具有如下形式
\begin{equation}
  \Pi^{\mu\nu ab}(k) = i\Bigl( g^{\mu\nu} - \frac{k^\mu k^\nu}{k^2}
    \Bigr)\ \bigl( m^2_{ab} + \mathcal{O}(k^2) \bigr) .
  \label{eq:20.53}
\end{equation}
在这种一般情形下计算 \eqref{eq:20.53} 中的非奇异项并不容易，但计算奇异项则是直接的，它来自具有中间戈德斯通玻色子的贡献。综合式~\eqref{eq:20.45} 和 \eqref{eq:20.46}，我们看到，规范玻色子转变为戈德斯通玻色子的振幅为
\begin{equation}
  \text{($A^a$-$\pi^j$ 顶角)} = -g k^\mu F^{aj} .
  \label{eq:20.54}
\end{equation}
于是对真空极化的极点贡献为
\begin{equation}
  \text{(极点贡献)} = \bigl( gk^\mu F^{aj} \bigr)\
    \frac{i}{k^2}\ \bigl( -gk^\nu F^{bj} \bigr) .
  \label{eq:20.55}
\end{equation}
把 \eqref{eq:20.55} 与式~(20.21) 比较，我们识别出
\begin{equation}
  m^2_{ab} = g^2\,F^{aj}F^{bj} .
  \label{eq:20.56}
\end{equation}
注意，在对称性被标量场破缺的情形中，这一结果回到 \eqref{eq:20.18}。然而，式~\eqref{eq:20.56} 适用于任何自发破缺对称性的理论，无论对称性破缺在拉格朗日量中是否显而易见，也无论它是否需要强相互作用或其他非微扰效应。于是，一个普遍结果是：任何耦合到自发破缺对称性之流的规范玻色子都获得质量。

\section{弱相互作用的 Glashow-Weinberg-Salam 理论}

由 Glashow、Weinberg 和 Salam（GWS）提出的弱相互作用与电磁相互作用的规范理论基于规范群 $SU(2)\times U(1)$。我们从弱同位旋 $SU(2)$ 和弱超荷 $U(1)$ 出发。$SU(2)$ 的规范玻色子是 $W^{1\mu}$、$W^{2\mu}$、$W^{3\mu}$，耦合常数为 $g$；$U(1)$ 的规范玻色子是 $B^\mu$，耦合常数为 $g'$。规范场的拉格朗日量为
\begin{equation}
  \mathcal{L}_{\mathrm{gauge}} = -\frac{1}{4}W^{a\mu\nu}W^a_{\mu\nu}
    - \frac{1}{4}B^{\mu\nu}B_{\mu\nu} ,
  \label{eq:20.58}
\end{equation}
其中
\begin{equation}
  W^{a\mu\nu} = \partial^\mu W^{a\nu} - \partial^\nu W^{a\mu}
    + g\epsilon^{abc}W^{b\mu}W^{c\nu},\qquad
  B^{\mu\nu} = \partial^\mu B^\nu - \partial^\nu B^\mu .
\end{equation}
作用在具有弱同位旋表示矩阵 $T^a$ 和弱超荷 $Y$ 的场上的协变导数为
\begin{equation}
  D_\mu = \partial_\mu - igW^a_\mu T^a - ig'B_\mu Y .
  \label{eq:20.65}
\end{equation}

对称性被复标量双重态自发破缺
\begin{equation}
  \phi = \begin{pmatrix}\phi^+\\ \phi^0\end{pmatrix}
\end{equation}
其超荷为 $Y = 1/2$。标量势为
\begin{equation}
  V(\phi) = \mu^2\phi^{\dagger}\phi + \lambda(\phi^{\dagger}\phi)^2 ,
\end{equation}
其中 $\mu^2 < 0$，因此 $\phi$ 获得真空期望值
\begin{equation}
  \langle\phi\rangle = \frac{1}{\sqrt{2}}\begin{pmatrix}0\\ v\end{pmatrix} .
  \label{eq:20.58b}
\end{equation}
把它代入标量动能项，我们得到规范玻色子质量项。带电规范玻色子
\begin{equation}
  W^{\pm\mu} = \frac{1}{\sqrt{2}}\bigl( W^{1\mu} \mp iW^{2\mu} \bigr)
  \label{eq:20.63}
\end{equation}
获得质量
\begin{equation}
  m_W = \frac{1}{2}gv .
\end{equation}
中性规范玻色子发生混合。光子和 $Z^0$ 是质量本征态
\begin{equation}
  A_\mu = B_\mu\cos\theta_W + W^3_\mu\sin\theta_W,\qquad
  Z_\mu = -B_\mu\sin\theta_W + W^3_\mu\cos\theta_W ,
  \label{eq:20.64}
\end{equation}
其中弱混合角 $\theta_W$ 由下式定义
\begin{equation}
  \sin\theta_W = \frac{g'}{\sqrt{g^2+g'^2}},\qquad
  \cos\theta_W = \frac{g}{\sqrt{g^2+g'^2}} .
\end{equation}
光子保持无质量，$Z^0$ 获得质量
\begin{equation}
  m_Z = \frac{1}{2}\sqrt{g^2+g'^2}\,v = \frac{m_W}{\cos\theta_W} .
  \label{eq:20.73}
\end{equation}

我们可以把协变导数 \eqref{eq:20.65} 改写为如下形式
\begin{equation}
  D_\mu = \partial_\mu - \frac{ig}{\sqrt{2}}\bigl( W^{+\mu}T^+ + W^{-\mu}T^-
    \bigr) - \frac{ig}{\cos\theta_W}\ Z^\mu\bigl( T^3 - Q\sin^2\theta_W
      \bigr) - ieQA_\mu ,
  \label{eq:20.71}
\end{equation}
其中
\begin{equation}
  e = g\sin\theta_W ,
  \label{eq:20.72}
\end{equation}
而 $Q = T^3 + Y$ 是电荷。我们在此看到，所有弱玻色子的耦合都由两个参数描述：已被精确测量的电子电荷 $e$，以及一个新参数 $\theta_W$。由 $W$ 和 $Z$ 交换所诱导的耦合还将涉及这些粒子的质量。然而，这些质量并不独立，因为由式~\eqref{eq:20.63} 可得
\begin{equation}
  m_W = m_Z\cos\theta_W .
  \label{eq:20.73b}
\end{equation}
$W$ 和 $Z$ 交换过程的所有效应，至少在树图级别，都可以用三个基本参数 $e$、$\theta_W$ 和 $m_W$ 来写出。

\subsection{与费米子的耦合}

一旦费米子场的量子数被指定，协变导数 \eqref{eq:20.71} 就唯一地确定 $W$ 和 $Z^0$ 场与费米子的耦合。为确定这些量子数，我们必须考虑第~17.3 节中提到的事实：$W$ 玻色子只耦合到夸克和轻子的左手螺旋态。

在经典拉格朗日量的层面，构造费米子场的左手与右手分量以不同方式耦合到规范玻色子的理论并无困难。%
\footnote{在第~19.4 节中，我们论证过，这一策略在量子修正层面可能存在一个问题。我们将在下面检查我们构造的具体模型是否避免这一问题。} 早在第~3.2 节我们就看到，Dirac 费米子的动能项分裂为左手场和右手场各自独立的部分：
\begin{equation}
  \bar{\psi}i\slashed{\partial}\psi = \bar{\psi}_L i\slashed{\partial}\psi_L
    + \bar{\psi}_R i\slashed{\partial}\psi_R .
  \label{eq:20.74}
\end{equation}
当我们把 $\psi$ 耦合到规范场时，我们可以把 $\psi_L$ 和 $\psi_R$ 指派给规范群的不同表示。于是 \eqref{eq:20.74} 右端的两个项将含有两个不同的协变导数，而这些将意味着两套不同的耦合。

在 GWS 模型中，我们可以利用这一技术来确保只有夸克场和轻子场的左手分量耦合到 $W$ 玻色子。我们把左手费米子场指派为 $SU(2)$ 的双重态，而使右手费米子场在该群下为单态。一旦我们为每个费米子场指定了 $T^3$ 值，我们必须指派的 $Y$ 值就由 $Q = T^3 + Y$ 决定。这意味着 $Y$ 的指派对夸克和轻子的左手与右手分量也将不同。对于右手场，$T^3 = 0$，于是我们通过把 $Y$ 指派为等于电荷来重现标准电荷。例如，对右手 $u$ 夸克场，$Y = +2/3$；对 $e_R$，$Y = -1$。对左手场，
\begin{equation}
  E_L = \begin{pmatrix}\nu_L\\ e_L\end{pmatrix},\qquad
  Q_L = \begin{pmatrix}u_L\\ d_L\end{pmatrix} ,
  \label{eq:20.75}
\end{equation}
指派 $Y = -1/2$ 和 $Y = +1/6$（分别）与 $T^3 = \pm 1/2$ 结合，给出正确的电荷指派。由于左手与右手费米子处于基本规范群的不同表示中，把这两个分量视为不同的粒子、并由费米子质量项把它们混合起来，常常是有用的。

事实上，费米子质量项的构造是一个严重的问题，因为所有可能的此类项都被全局规范不变性所禁止。例如，我们不能写出电子质量项
\begin{equation}
  \Delta\mathcal{L} = -m_e\bigl( \bar{e}_Le_R + \bar{e}_Re_L \bigr) ,
  \label{eq:20.76}
\end{equation}
因为场 $e_L$ 和 $e_R$ 属于不同的 $SU(2)$ 表示，并且具有不同的 $U(1)$ 荷。在接下来的几页中，我们将通过把所有费米子场视为无质量来忽略这一问题。这一描述足以分析高能下弱相互作用的结构，在那种情况下夸克和轻子质量可以被忽略。在本节末尾，我们将回到在 GWS 理论中写出夸克和轻子质量项的问题。这一问题的解将强化这样一种观念：左手与右手费米子场根本上是独立的实体，由某种辅助过程混合形成有质量费米子。

如果我们忽略费米子质量，夸克和轻子弱相互作用的拉格朗日量就直接来自上面给出的电荷指派。$e$、$\nu$、$u$ 和 $d$ 的费米子动能项为
\begin{equation}
  \mathcal{L} = \bar{E}_L(i\slashed{D})E_L + \bar{e}_R(i\slashed{D})e_R
    + \bar{Q}_L(i\slashed{D})Q_L + \bar{u}_R(i\slashed{D})u_R
    + \bar{d}_R(i\slashed{D})d_R .
  \label{eq:20.77}
\end{equation}
在每一项中，协变导数由式~\eqref{eq:20.65} 给出，其中 $T^a$ 和 $Y$ 在所属费米子场的特定表示中取值。例如，
\begin{equation}
  \bar{Q}_L(i\slashed{D})Q_L = \bar{Q}_L i\gamma^\mu\Bigl( \partial_\mu
    - igW^a_\mu T^a - i\frac{g'}{6}B_\mu \Bigr)Q_L .
  \label{eq:20.78}
\end{equation}
右手中微子对 $SU(2)$ 和 $U(1)$ 的耦合都为零，因此我们只是从式~\eqref{eq:20.77} 中省略了这一场。

为弄清费米子-矢量玻色子耦合的物理后果，我们应该利用式~\eqref{eq:20.71} 所给协变导数的形式，把式~\eqref{eq:20.77} 用矢量玻色子质量本征态写出。于是式 \eqref{eq:20.77} 具有如下形式
\begin{equation}
  \begin{aligned}
    \mathcal{L} &= \bar{E}_L(i\slashed{\partial})E_L + \bar{e}_R(i\slashed{\partial})e_R
      + \bar{Q}_L(i\slashed{\partial})Q_L + \bar{u}_R(i\slashed{\partial})u_R
      + \bar{d}_R(i\slashed{\partial})d_R\\
    &\quad + g\bigl( W^+_\mu J^{\mu+} + W^-_\mu J^{\mu-} + Z^0_\mu J^\mu_Z
      \bigr) + eA_\mu J^\mu_{EM} ,
  \end{aligned}
  \label{eq:20.79}
\end{equation}
其中
\begin{equation}
  \begin{aligned}
    J^{\mu+}_W &= \frac{1}{\sqrt{2}}\bigl( \bar{\nu}_L\gamma^\mu e_L
      + \bar{u}_L\gamma^\mu d_L \bigr),\\
    J^{\mu-}_W &= \frac{1}{\sqrt{2}}\bigl( \bar{e}_L\gamma^\mu\nu_L
      + \bar{d}_L\gamma^\mu u_L \bigr),\\
    J^\mu_Z &= \frac{1}{\cos\theta_W}\ \Bigl[ \bar{\nu}_L\gamma^\mu
      \Bigl( \frac{1}{2} \Bigr)\nu_L + \bar{e}_L\gamma^\mu\Bigl( -\frac{1}{2}
        + \sin^2\theta_W \Bigr)e_L + \bar{e}_R\gamma^\mu(\sin^2\theta_W)e_R\\
    &\quad + \bar{u}_L\gamma^\mu\Bigl( \frac{1}{2} - \frac{2}{3}\sin^2\theta_W
      \Bigr)u_L + \bar{u}_R\gamma^\mu\Bigl( -\frac{2}{3}\sin^2\theta_W
        \Bigr)u_R\\
    &\quad + \bar{d}_L\gamma^\mu\Bigl( -\frac{1}{2} + \frac{1}{3}\sin^2\theta_W
      \Bigr)d_L + \bar{d}_R\gamma^\mu\Bigl( \frac{1}{3}\sin^2\theta_W
        \Bigr)d_R \Bigr],\\
    J^\mu_{EM} &= \bar{e}\gamma^\mu(-1)e + \bar{u}\gamma^\mu\Bigl(
      +\frac{2}{3} \Bigr)u + \bar{d}\gamma^\mu\Bigl( -\frac{1}{3} \Bigr)d .
  \end{aligned}
  \label{eq:20.80}
\end{equation}
在此我们使用了式~(20.67) 来简化 $W$ 玻色子流。注意与光子场相关联的流 $J^\mu_{EM}$ 确实就是标准的电磁流。

\subsection{反常消除}

正如我们刚才所见，写出一个以手征方式把 GWS 规范玻色子耦合到费米子的拉格朗日量并无困难。然而，这些手征耦合确实呈现一个在单圈修正层面出现的问题。在第~19.2 节中，我们看到，在经典运动方程层面守恒的轴矢流可以通过把该流耦合到一对规范玻色子的单圈图而获得非零散度。含有这一反常贡献的费曼图是一个三角图，其顶点处为轴矢流和两个规范流。在规范玻色子耦合到手征流的规范理论中，危险的三角图出现在三规范玻色子顶角函数的单圈修正中。反常项破坏了这一振幅的 Ward 恒等式。因此，如我们在第~19.4 节所论证的，规范玻色子耦合到手征流的理论只有在反常贡献以某种方式消失时才能是规范不变的。幸运的是，如我们在那里所见，当人们对可以在这些图中循环的所有可能的费米子种类求和时，反常项可以安排得相互抵消。%
\footnote{如果你没有读过第~19 章，但你愿意假定费米子三角图含有破坏规范不变性的贡献，你仍然应该能够跟上接下来的论证。}

在 GWS 理论中，实验要求弱相互作用流是左手的，这迫使我们选择手征规范耦合。现在我们必须检查来自三角图的反常项是否按要求抵消。我们将发现它们确实如此，但只是通过夸克和轻子量子数之间微妙而相当奇妙的相互作用才实现。

三个规范玻色子 $A^a$、$A^b$ 和 $A^c$ 的三角图的反常项正比于群论不变量
\begin{equation}
  \mathrm{tr}\Bigl[ \gamma^5\ t^a\{t^b, t^c\} \Bigr] ,
  \label{eq:20.81}
\end{equation}
其中求迹遍及所有费米子种类。反对易子来自取费米子沿相反方向绕行的两个三角图之和。因子 $\gamma^5$ 记录反常与手征流相关联的事实；该因子对左手费米子等于 $-1$，对右手费米子等于 $+1$。在诸如 QED 或 QCD 这类规范玻色子对右手与左手种类等权耦合的理论中，反常自动抵消。这种记账方法是第~19.4 节所给更一般方法的一个特例。

为求值 GWS 理论的反常，最容易的做法是在 $SU(2)\times U(1)$ 规范玻色子的基中、在混合到光子和 $Z^0$ 质量本征态之前进行。对无质量费米子求值三角图就足够了，这样右手与左手费米子具有不同的量子数。然而，我们不仅要考虑含三个 $SU(2)\times U(1)$ 规范玻色子的图的反常，还要考虑既含弱相互作用规范玻色子又含 QCD 色 $SU(3)$ 规范玻色子的图。如果我们考虑引力对弱相互作用规范理论的影响，还存在一个含一个弱相互作用规范玻色子和两个引力子的可能反常的图。我们可以忽略所有耦合都是左右对称的图，例如三个 $SU(3)$ 玻色子的反常或一个 $SU(3)$ 玻色子与两个引力子的反常。于是，可能含有反常项的完整图集如图~\ref{fig:20.2} 所示。如果 $SU(2)\times U(1)$ 规范理论的 Ward 恒等式要得到满足，所有可能的反常都必须消除。

\begin{figure}[t]
  \centering
  \includegraphics[width=0.6\textwidth]{fig20.2.png}
  \caption{在 GWS 理论中可能产生反常的三角图。}
  \label{fig:20.2}
\end{figure}

三个 $SU(2)$ 规范玻色子的反常总为零是 $SU(2)$ 规范理论的一个特殊性质；这一结果来自 Pauli 西格马矩阵的性质 $\{\sigma^a,\sigma^b\} = 2\delta^{ab}$，它意味着迹 \eqref{eq:20.81} 为零。含一个 $SU(3)$ 玻色子或一个 $SU(2)$ 玻色子的反常正比于
\begin{equation}
  \mathrm{tr}[t^a] = 0
  \qquad\text{或}\qquad
  \mathrm{tr}[\tau^a] = 0 .
  \label{eq:20.82}
\end{equation}
其余非平凡的反常是一个 $U(1)$ 玻色子与两个 $SU(2)$ 玻色子或两个 $SU(3)$ 玻色子的反常、三个 $U(1)$ 玻色子的反常，以及含一个 $U(1)$ 规范玻色子的引力反常。

一个 $U(1)$ 玻色子与两个 $SU(3)$ 玻色子的反常正比于群论因子
\begin{equation}
  \mathrm{tr}\bigl[ t^at^bY \bigr] = \frac{1}{2}\delta^{ab}\cdot
    \sum_{\text{夸克}} Y ,
  \label{eq:20.83}
\end{equation}
其中求和遍及左手夸克和右手夸克，对左手贡献多一个因子 $(-1)$。代入上面为 $u_L$、$d_L$、$u_R$ 和 $d_R$ 给出的电荷指派，我们得到
\begin{equation}
  \sum_{\text{夸克}} Y = -2\cdot\frac{1}{6} + \Bigl( \frac{2}{3} \Bigr)
    + \Bigl( -\frac{1}{3} \Bigr) = 0 .
  \label{eq:20.84}
\end{equation}
一个 $U(1)$ 玻色子与两个 $SU(2)$ 玻色子的反常正比于左手双重态超荷之和。把轻子双重态和夸克双重态加起来，我们得到
\begin{equation}
  \sum_{\text{双重态}} Y = 3\cdot\Bigl( \frac{1}{6} \Bigr)
    + \Bigl( -\frac{1}{2} \Bigr) = 0 .
\end{equation}
三个 $U(1)$ 玻色子的反常正比于超荷立方之和。把所有左手与右手费米子的贡献按适当的符号加起来，人们发现这也为零，其中夸克带有色因子 3：
\begin{equation}
  \sum_{\text{所有费米子}} Y^3 = 3\cdot\Bigl[ -2\Bigl( \frac{1}{6} \Bigr)^3
    + \Bigl( \frac{2}{3} \Bigr)^3 + \Bigl( -\frac{1}{3} \Bigr)^3 \Bigr]
    - 2\Bigl( -\frac{1}{2} \Bigr)^3 + (-1)^3 = 0 .
  \label{eq:20.85}
\end{equation}
最后，含一个 $U(1)$ 规范玻色子的引力反常正比于超荷之和，我们已证明它为零。因此，GWS 理论是无反常的。

GWS 理论中反常的消除是一个显著而非平凡的自洽性条件。它要求夸克具有分数电荷，轻子具有整数电荷，并且颜色数为 3。这些条件都被观测到的粒子谱所满足，这一事实有力地表明 GWS 理论走在正确的道路上。

\subsection{Glashow-Weinberg-Salam 理论的实验后果}

GWS 理论对弱相互作用现象作出若干精确预言。在低能下，重的 $W$ 和 $Z$ 玻色子的交换产生有效的四费米子相互作用。带电流相互作用重现了弱相互作用 Fermi 理论的熟悉结构，其中 Fermi 常数由下式给出
\begin{equation}
  \frac{G_F}{\sqrt{2}} = \frac{g^2}{8m_W^2} .
  \label{eq:20.91}
\end{equation}
中性流相互作用给出 Fermi 理论中不存在的新的预言。这些预言已经在各种各样的实验中受到检验。

在更高能量下还有对 GWS 理论的进一步检验。过程 $e^{+}e^{-}\to f\bar{f}$ 以本质性的方式受到影响，因为该理论含有一个带 $s$ 道 $Z^0$ 交换的新图，它与标准的光子交换图干涉，如图~\ref{fig:20.4} 所示。利用第~5.2 节的方法计算出这一干涉的效应是直接的，因此我们把这一分析留作习题~20.3。

\begin{figure}[t]
  \centering
  \includegraphics[width=0.55\textwidth]{fig20.4.png}
  \caption{在 Glashow-Weinberg-Salam 理论中对过程 $e^{+}e^{-}\to f\bar{f}$ 有贡献的图。}
  \label{fig:20.4}
\end{figure}

当质心能量接近 $m_Z$ 时，$Z^0$ 直接作为 $e^{+}e^{-}$ 湮灭截面中的共振峰出现。类似地，$W$ 和 $Z$ 都可以在夸克-反夸克湮灭中作为共振峰被观测到，后者被视为质子-反质子散射中的部分子次级过程。这些共振峰的位置由 $G_F$、$\sin^2\theta_W$ 以及 $e$ 或 $\alpha$ 的值根据式~\eqref{eq:20.72} 和 \eqref{eq:20.91} 预言。利用这些关系，我们得到
\begin{equation}
  m_W = \sqrt{\frac{\alpha\pi}{\sqrt{2}G_F\sin^2\theta_W}},\qquad
  m_Z = \sqrt{\frac{\alpha\pi}{\sqrt{2}G_F\sin^2\theta_W\cos^2\theta_W}} .
  \label{eq:20.95}
\end{equation}
$Z^0$ 共振峰的详细形状如图~\ref{fig:20.5} 所示。所示的实验测量结果与一条共振峰位置按最佳拟合调整过的理论曲线相比较。共振峰的高度和宽度随后由 GWS 理论预言。共振峰通过电子和正电子在湮灭前辐射共线光子的过程而被展宽到更高能量；这一修正在习题~5.5 中讨论过。

\begin{figure}[t]
  \centering
  \includegraphics[width=0.6\textwidth]{fig20.5.png}
  \caption{在 $E_{cm}$ 接近 $Z^0$ 玻色子质量时，$e^{+}e^{-}$ 湮灭为强子的总截面，由 ALEPH、DELPHI、L3 和 OPAL 实验测量，并由粒子数据组（Particle Data Group）汇编，\emph{Phys.\ Rev.\ D} \textbf{50}, (1994), 图~32.14。原始文章的参考文献列于其中。实线是 GWS 理论的预言。}
  \label{fig:20.5}
\end{figure}

由于 GWS 理论的拉格朗日量把左手与右手费米子视为具有完全不同量子数的不同种类，$Z^0$ 对左手与右手费米子的耦合显著不同。这一点的表现之一是存在极化不对称性，即衰变 $Z^0\to f\bar{f}$ 中产生的费米子具有净极化，或者在逆过程即 $Z^0$ 产生中存在不对称性。这一不对称性可以直接从 \eqref{eq:20.80} 所给 $Z^0$ 流的形式读出：
\begin{equation}
  \mathcal{A}_{LR} = \frac{\Gamma(Z^0\to f_L\bar{f}_R)
    - \Gamma(Z^0\to f_R\bar{f}_L)}
    {\Gamma(Z^0\to f_L\bar{f}_R) + \Gamma(Z^0\to f_R\bar{f}_L)}
  = \frac{\bigl( \tfrac{1}{2} - |Q_f|\sin^2\theta_W \bigr)^2
    - \bigl( Q_f\sin^2\theta_W \bigr)^2}
    {\bigl( \tfrac{1}{2} - |Q_f|\sin^2\theta_W \bigr)^2
    + \bigl( Q_f\sin^2\theta_W \bigr)^2} .
  \label{eq:20.96}
\end{equation}
对现实的取值 $\sin^2\theta_W = 0.23$，该表达式对带电轻子给出 15\% 的不对称性，对 $d$、$s$ 和 $b$ 夸克给出 95\% 的不对称性。对轻子，这一不对称性可以通过在 $Z^0$ 共振峰处测量 $\tau$ 轻子的极化、或者通过测量分别用左手与右手电子产生共振峰的相对截面而在实验上加以检验。对夸克，该不对称性可以通过共振峰上的前后产生不对称性间接确定，如习题~20.3 所解释。

由于弱中性流有如此多不同的表现，通过比较为解释其每一个预言效应所需的参数 $\sin^2\theta_W$ 的值，GWS 弱相互作用理论可以受到严格的检验。表~\ref{tab:20.1} 给出从各种各样的弱相互作用中性流效应和不对称性中提取的 $\sin^2\theta_W$ 值。在所有情形下，都必须包含单圈辐射修正才能在所要求的精度水平上分析实验。这些辐射修正涉及一些微妙之处。首先，必须采用一个定义 $\sin^2\theta_W$ 的特定重整化约定，并在所有计算中一致地使用它。该表给出这一约定的两种不同选择的结果。在这两种约定中，弱相互作用可观测量都被取为 $\alpha$、$G_F$ 和第三个独立参数的函数。

\begin{table}[t]
  \centering
  \caption{弱相互作用实验给出的 $\sin^2\theta_W$ 值}
  \begin{tabular}{lll}
    \toprule
    $\sin^2\theta^{MS}_W$ & $\sin^2\theta^{MSbar}_W$ & 观测量或过程\\
    \midrule
    0.2247 (21) & 0.2320 (6) & $m_Z$\\
    0.2264 (25) & 0.2338 (22) & $m_W$\\
    0.2250 (18) & 0.2322 (6) & $\Gamma_Z$\\
    0.2243 (17) & 0.2315 (11) & 在 $Z^0$ 处轻子的 $f$-$b$ 不对称性\\
    0.2245 (17) & 0.2317 (8) & 在 $Z^0$ 处所有对产生不对称性\\
    0.2221 (17) & 0.2292 (10) & 在 $Z^0$ 处的 $A_{LR}$\\
    0.2260 (48) & 0.233 (5) & 深非弹中微子散射\\
    0.205 (31) & 0.212 (32) & 中微子-质子弹性散射\\
    0.224 (9) & 0.231 (9) & 中微子-电子弹性散射\\
    0.216 (8) & 0.223 (8) & 原子宇称破坏\\
    0.216 (17) & 0.223 (18) & 非弹 $e^{-}$ 散射中的宇称破坏\\
    \bottomrule
  \end{tabular}
  \label{tab:20.1}
\end{table}

此处列出的值是通过调整 $\sin^2\theta^{MS}_W$ 或 $\sin^2\theta^{\overline{MS}}_W$ 的值、把 $\alpha$ 和 $G_F$ 视为精确已知的参数来拟合实验观测而得到的。括号中的数字是末位所示数字的标准误差。从实验测量量到 $\sin^2\theta_W$ 或 $\sin^2\theta^{\overline{MS}}_W$ 的换算依赖于顶夸克质量和希格斯玻色子质量。这些值假设顶夸克质量为 169 GeV、希格斯质量为 300 GeV；所引用的误差包括顶夸克质量的 17 GeV 不确定度以及希格斯质量从 60 GeV 到 1000 GeV 的范围。两列之间相对误差的差异反映了这一理论不确定性的重要性。一些可观测量对 $\alpha_s$ 依赖较弱；这些值假设 $\alpha_s(m_Z) = 0.120 \pm 0.007$。该表取自 P.~Langacker 和 J.~Erler 为粒子数据组（Particle Data Group）撰写的文章，\emph{Phys.\ Rev.\ D} \textbf{50}，1304 (1994)。那篇文章包含完整的参考文献集，并讨论了这些确定中不确定性的来源。

在第一列中，该参数是质量比 $m_W/m_Z$，并且，遵循树图级表达式 \eqref{eq:20.73}，我们认为这一比值定义了 $\sin^2\theta_W$ 的重整化值：
\begin{equation}
  \sin^2\theta^s_W = 1 - \frac{m_W^2}{m_Z^2} .
  \label{eq:20.97}
\end{equation}
在第二列中，第三个参数是根据由最小减除（式~\eqref{eq:11.77}）定义的弱相互作用耦合常数计算出的 $\sin^2\theta_W$。$\sin^2\theta_W$ 不同定义之间的差异出现在单圈计算的层面，并且可以揭示有趣的物理；这一主题在第~21.3 节讨论。

第二个微妙之处在于，弱中性流过程的单圈修正依赖于 $t$ 夸克质量的值，这一质量直到最近才被确定，并且仍然有些认识不足。对 $t$ 夸克质量的依赖相对较强，其原因很有意思，我们将在第~21.3 节讨论。单圈修正还对负责自发对称性破缺的粒子的性质有较弱的依赖。

从表~\ref{tab:20.1} 我们可以看到，由弱中性流引起的各种各样的效应都已被观测到，其大小由单一且自洽的 $\sin^2\theta_W$ 值所解释。理论与实验的这种显著一致使我们相信，Glashow-Weinberg-Salam 理论确实是弱相互作用与电磁相互作用的正确描述。

\subsection{费米子质量项}

我们现在回到为夸克和轻子写出质量项的问题。回忆我们不能把普通质量项放入拉格朗日量中，因为各种费米子场的左手与右手分量具有不同的规范量子数，所以简单的质量项会破坏规范不变性。为赋予夸克和轻子质量，我们必须再次援引自发对称性破缺的机制。

我们在本节开头假定标量场 $\phi$ 获得真空期望值 \eqref{eq:20.58}，以便赋予 $W$ 和 $Z$ 玻色子质量。该标量场在 $SU(2)$ 下必须是旋量，并且必须具有 $Y = 1/2$，才能产生正确的规范玻色子质量模式。利用这些量子数，我们还可以写出把 $e_L$、$e_R$ 和 $\phi$ 联系起来的规范不变耦合，如下所示：
\begin{equation}
  \Delta\mathcal{L}_e = -\lambda_e\ \bar{E}_L\cdot\phi\,e_R + \mathrm{h.c.}
  \label{eq:20.98}
\end{equation}
在此，双重态 $E_L$ 和 $\phi$ 的 $SU(2)$ 指标被缩并；还要注意各个场的荷 $Y$ 之和为零。参数 $\lambda_e$ 是一个新的无量纲耦合常数。如果我们把该表达式中的 $\phi$ 换成其真空期望值 \eqref{eq:20.58}，就得到
\begin{equation}
  \Delta\mathcal{L}_e = -\frac{1}{\sqrt{2}}\lambda_e\,v\ \bar{e}_Le_R
    + \mathrm{h.c.} + \cdots .
  \label{eq:20.99}
\end{equation}
这是电子的质量项。质量的大小由 $\phi$ 的真空期望值设定，并被新的无量纲耦合重新标度：
\begin{equation}
  m_e = \frac{1}{\sqrt{2}}\lambda_e v .
  \label{eq:20.100}
\end{equation}
由于电子质量正比于 $v$，人们可能会预期电子和 $W$ 玻色子的质量应当同数量级。事实上，取观测值，$m_e/m_W \sim 6\times10^{-6}$。由于 $\lambda_e$ 是可重整化耦合，它必须被当作理论的输入。因此，GWS 理论允许电子非常轻，但它不能解释为什么电子相对于 $W$ 玻色子如此之轻。

我们可以用同样的方式写出夸克场的质量项。注意，在下面的表达式中，两项都在 $SU(2)$ 下不变，并且净 $Y$ 为零：
\begin{equation}
  \Delta\mathcal{L}_q = -\lambda_d\ \bar{Q}_L\cdot\phi\,d_R
    - \lambda_u\ \epsilon_{ab}\ \bar{Q}^a_L\phi^{b*}\,u_R + \mathrm{h.c.}
  \label{eq:20.101}
\end{equation}
代入来自式~\eqref{eq:20.58} 的 $\phi$ 的真空期望值，这些项变为
\begin{equation}
  \Delta\mathcal{L}_q = -\frac{1}{\sqrt{2}}\lambda_d\,v\ \bar{d}_Ld_R
    - \frac{1}{\sqrt{2}}\lambda_u\,v\ \bar{u}_Lu_R + \mathrm{h.c.} + \cdots ,
  \label{eq:20.102}
\end{equation}
即 $d$ 和 $u$ 夸克的标准质量项。GWS 理论由此给出关系
\begin{equation}
  m_d = \frac{1}{\sqrt{2}}\lambda_d v,\qquad
  m_u = \frac{1}{\sqrt{2}}\lambda_u v .
  \label{eq:20.103}
\end{equation}
与电子一样，该理论参数化但不解释实验上观测到的 $d$ 和 $u$ 夸克质量的小值。

当在理论中引入额外几代夸克时，可以存在混合世代的额外耦合项。或者，我们可以通过为夸克场选择新基来对角化 Higgs 耦合。我们将在第~20.3 节证明这总是可能的。然而，Higgs 耦合的这一简化会引起规范耦合中的复杂化。于是带电流弱相互作用通过一个幺正矩阵——Cabibbo-Kobayashi-Maskawa（CKM）矩阵——耦合各代夸克，该矩阵一般不是对角的。

希格斯玻色子的唯象性质将在第~III 部分的总课题中更详细地探讨。

\subsection{一个希格斯扇区？}

由于不存在支持 GWS 模型中所含简单希格斯玻色子存在的实验证据，值得追问 $W$ 和 $Z$ 玻色子是否可能通过更复杂的机制获得质量。这一问题有两个方面。

首先，$W$ 和 $Z$ 玻色子是否确实是自发破缺的 $SU(2)\times U(1)$ 对称性的规范玻色子，这确定无疑吗？支持这一想法的证据来自各种夸克和轻子对弱相互作用的耦合的普适性。这一普适性通过如下事实受到检验：同一个 Fermi 常数值描述了所有带电流弱相互作用过程，而同样的耦合强度与单一的 $\sin^2\theta_W$ 值相结合描述了整个弱中性流现象范围。我们已经看到，尤其是在第~16 章的讨论中，局域规范不变性原理自然而然地导致普适的、与味无关的耦合常数的预言。人们不知道还有任何其他原理能解释这一引人注目的规律性。因此，有令人信服的证据表明，弱相互作用的底层理论是自发破缺的规范理论。

然而，$SU(2)\times U(1)$ 的自发破缺机制很可能比我们在式~(20.111) 中写下的单一标量场的简单模型更复杂。原则上，$SU(2)\times U(1)$ 的破缺可能是由一组复杂的新粒子和新相互作用的动力学造成的，我们称之为\emph{希格斯扇区}。实验只给出这一新扇区的三个性质：第一，它必须产生夸克和轻子的质量。第二，它必须产生 $W$ 和 $Z$ 玻色子的质量。第三项信息——也是唯一非平凡的一项——来自 GWS 理论中弱玻色子质量之间的关系 \eqref{eq:20.73}，
\begin{equation}
  m_W = m_Z\cos\theta_W .
  \label{eq:20.117}
\end{equation}
这一关系在实验上以好于 1\% 的精度成立，也就是说，达到了单圈辐射修正的水平。无论我们援引什么复杂的机制来产生 $SU(2)\times U(1)$ 的自发破缺，它都应当以自然的方式重现这一关系。

为理解关系 \eqref{eq:20.117} 的含义，我们必须在不假设 $SU(2)\times U(1)$ 被标量场的期望值破缺的情况下分析规范玻色子质量矩阵。实际上，在宽松得多的假设下，利用第~20.1 节最末尾给出的论证，就有可能计算出规范玻色子质量矩阵。在那里，我们从规范流产生或湮灭戈德斯通玻色子的矩阵元构造了规范玻色子质量矩阵。我们现在将证明，对一大类这些矩阵元满足某些简单性质的 $SU(2)\times U(1)$ 破缺模型，关系 \eqref{eq:20.117} 都成立。

任何弱相互作用对称性破缺模型都必须含有某组负责 $SU(2)\times U(1)$ 自发破缺的场。把理论的这一扇区视为具有全局 $SU(2)\times U(1)$ 对称性的场论，该对称性通过其与规范玻色子的耦合被提升为局域对称性。在具有全局对称性的理论中，该对称性被自发破缺为 $U(1)$。由于三个连续对称性被自发破缺，该扇区必须提供三个戈德斯通玻色子，它们最终将被 $W^+$、$W^-$ 和 $Z^0$ 吃掉。称这三个玻色子为 $\pi^a$，其中 $a = 1, 2, 3$。设 $J^{\mu A}$ 是新扇区的 $SU(2)$ 对称性流，并设 $J^{\mu Y}$ 是 $U(1)$ 流。于是规范玻色子质量矩阵将由矩阵元 \eqref{eq:20.46} 构造，后者在此具有如下形式
\begin{equation}
  \langle 0|J^{\mu A}|\pi^b(p)\rangle = -ip^\mu F^{Ab} ,
  \label{eq:20.118}
\end{equation}
其中 $A = 1, 2, 3, Y$，$b = 1, 2, 3$。利用式~\eqref{eq:20.55} 的方法，我们发现规范玻色子真空极化含有极点项
\begin{equation}
  -\frac{i}{k^2}\ \bigl( g_AF^{Ac} \bigr)\bigl( g_BF^{Bc} \bigr) ,
  \label{eq:20.119}
\end{equation}
其中对 $c$ 求和，$A = 1, 2, 3$ 时 $g_A = g$，$A = Y$ 时 $g_A = g'$。于是我们可以把规范玻色子质量矩阵识别为
\begin{equation}
  m^2_{AB} = g_Ag_B\,F^{Ac}F^{Bc} .
  \label{eq:20.120}
\end{equation}
为重现已知形式的弱规范玻色子质量矩阵，我们现在必须对 $F^{Ab}$ 施加约束。首先我们必须确保光子保持无质量。如果电荷的线性组合式~(20.69) 湮灭真空，这一点就成立。用式~\eqref{eq:20.118} 的语言来说，我们必须坚持相应的电流线性组合不能激发戈德斯通玻色子：
\begin{equation}
  \langle 0|\bigl( J^{\mu 3} + J^{\mu Y} \bigr)|\pi^0(p)\rangle = 0 .
  \label{eq:20.121}
\end{equation}
利用下面的附加假设，我们也能实现关系 \eqref{eq:20.117}：对称性破缺扇区具有 $SU(2)$ 全局对称性，在该对称性下三个戈德斯通玻色子和三个 $SU(2)$ 规范流作为三重态变换，并且当 $SU(2)$ 规范对称性被自发破缺时它仍然保持精确。这一全局 $SU(2)$ 对称性意味着，如果在式~\eqref{eq:20.118} 中 $A = a = 1, 2, 3$，那么
\begin{equation}
  \langle 0|J^{\mu a}|\pi^b(p)\rangle = -iF\,p^\mu\,\delta^{ab} ,
  \label{eq:20.122}
\end{equation}
其中 $F$ 是具有质量量纲的参数。综合 \eqref{eq:20.122} 和 \eqref{eq:20.121}，我们得到
\begin{equation}
  \langle 0|J^{\mu Y}|\pi^3(p)\rangle = +iF\,p^\mu .
  \label{eq:20.123}
\end{equation}
把 $F^{Ab}$ 的这一形式代入 \eqref{eq:20.120}，我们得到规范玻色子质量矩阵
\begin{equation}
  m^2 = F^2\ \begin{pmatrix}
    g^2 & & & \\
    & g^2 & & \\
    & & g^2 & -gg'\\
    & & -gg' & g'^2
  \end{pmatrix},
  \label{eq:20.124}
\end{equation}
其中矩阵作用在规范玻色子 $(A^1_\mu, A^2_\mu, A^3_\mu, B_\mu)$ 上。该矩阵的特征向量正是 \eqref{eq:20.63} 和 \eqref{eq:20.64}。为重现特征值，我们只需要定义
\begin{equation}
  v = 2F .
  \label{eq:20.125}
\end{equation}
我们现在已经证明，$W$ 和 $Z$ 玻色子质量之间的 GWS 关系并非规范对称性被单一标量场破缺的情形所特有。这一关系来自对希格斯扇区未破缺全局 $SU(2)$ 对称性的更普遍假设。该对称性常被称为\emph{守衡 $SU(2)$ 对称性}。%
\footnote{P.~Sikivie、L.~Susskind、M.~Voloshin 和 V.~Zakharov，\emph{Nucl.\ Phys.\ B} \textbf{173}，189 (1980)。} 我们已经在弱相互作用有效拉格朗日量式~(20.94) 中把它作为全局 $SU(2)$ 对称性看到过。

对单一标量场的情形，守衡对称性以下列方式出现：如果我们把场 $\phi$ 用其四个实分量写出，拉格朗日量式~(20.111)（忽略规范耦合）具有 $O(4)$ 全局对称性。$\phi$ 的真空期望值把该对称性破缺为 $O(3)$，即 $SU(2)$。

然而，还有许多其他量子场论在自发破缺 $SU(2)$ 的同时保持另一个全局 $SU(2)$ 对称性未破缺。一个相当复杂的例子由具有两个无质量味的 QCD 给出，如果我们把被规范的 $SU(2)$ 与 \eqref{eq:19.82} 中由 $U_L$ 生成的对称性等同起来，并把守衡 $SU(2)$ 与矢量同位旋对称性等同起来。把熟悉的强相互作用复制一份，其质量标度大到足以给出 $F = 125$ GeV，就将是希格斯扇区一个完全可以接受的模型。（遗憾的是，在这一模型中不容易产生夸克和轻子的质量。）

关于希格斯扇区的性质以及 $SU(2)\times U(1)$ 破缺的显式机制的问题，很可能是基本粒子理论中最紧迫的未解问题。我们将在结语中进一步讨论这一问题。

\section{夸克与轻子理论的对称性}

把第~17 章描述的强相互作用理论与上一节描述的弱相互作用和电磁相互作用理论放在一起，我们现在已经构造了基本粒子相互作用的完整描述。研究这一理论的对称性，追问自然界底层描述的基本对称性可能是什么，是很有意思的。

在通向式~\eqref{eq:15.17} 的论证中，我们已经看到，规范理论的拉格朗日量受到可重整化与规范不变性条件的严格限制。在本节中，我们将构造与强、弱和电磁相互作用的 $SU(3)\times SU(2)\times U(1)$ 规范对称性相容的最一般的可重整化拉格朗日量。然后我们可以追问，为了赋予该理论我们在自然界中看到的那些全局对称性，我们还必须对该理论施加哪些额外的全局对称性。

作为第一步，我们将忽略希格斯标量场以及夸克、轻子和规范玻色子的质量项。于是夸克与轻子理论的拉格朗日量完全由规范不变性和可重整化性所确定。我们有
\begin{equation}
  \mathcal{L}_{G} = -\frac{1}{4}\sum_i F^{i\mu\nu}F^i_{\mu\nu}
    + \sum_J \bar{\psi}_J i\slashed{D}\psi_J ,
  \label{eq:20.126}
\end{equation}
其中指标 $i$ 遍及规范群的三个因子，指标 $J$ 遍及手征费米子的各个多重态。

原则上，我们可以向 \eqref{eq:20.126} 添加如下的赝标量纯规范算符：
\begin{equation}
  \Delta\mathcal{L}_{\theta} = \sum_i \theta_i\ \frac{g_i^2}{32\pi^2}\,
    \epsilon^{\mu\nu\rho\sigma}F^i_{\mu\nu}F^i_{\rho\sigma} .
  \label{eq:20.127}
\end{equation}
这些项显然在 $P$ 和 $T$ 下都是奇的。然而，我们在第~19.2 节末尾看到，通过在有效作用量中做变量替换，可以产生或抵消这种形式的项。例如，对右手电子场的变量替换
\begin{equation}
  e_R \to e^{i\alpha}e_R
  \label{eq:20.128}
\end{equation}
根据 \eqref{eq:19.78} 或 \eqref{eq:19.79}，产生对拉格朗日量的一个修正，该修正涉及 $U(1)$ 规范场场强中 $P$ 奇与 $T$ 奇的组合
\begin{equation}
  \Delta\mathcal{L} = \alpha\,\frac{g'^2}{32\pi^2}\,
    \epsilon^{\mu\nu\rho\sigma}B_{\mu\nu}B_{\rho\sigma} .
  \label{eq:20.129}
\end{equation}
\eqref{eq:20.129} 的系数与 \eqref{eq:19.79} 中相应的系数不同，因为我们只变换了电子场的右手手征分量。如果我们变换另一个超荷为 $Y$ 的费米子场，我们会发现类似的移动，其系数正比于 $Y^2$。如果这一新场耦合到 $SU(2)$ 或 $SU(3)$ 规范场，我们还会发现正比于那些场强的项。因此，我们可以通过做取 $\alpha = -\theta_1$ 的变量替换 \eqref{eq:20.128}，消去 \eqref{eq:20.127} 中涉及 $U(1)$ 场强的项。我们可以通过对三个费米子多重态（比如说 $e_R$、$E_L$ 和 $Q_L$）做适当的手征转动，消去 \eqref{eq:20.127} 中的全部三项。变量替换 \eqref{eq:20.128} 转动右手电子场，在宇称下不是对称的，并且事实上改变了宇称运算的定义。通过做这一变量替换，我们是在选择新的坐标，使得整个理论的 $P$ 和 $T$ 变换性质尽可能简单。

现在让我们研究拉格朗日量 \eqref{eq:20.126} 在 $P$、$C$ 和 $T$ 下的性质。QCD 规范玻色子的耦合对这些对称性中的每一个都分别不变。然而，$SU(2)$ 规范玻色子的耦合尽可能地破坏 $P$ 和 $C$。回忆第~3.6 节，$P$ 把左手电子变为右手电子，$C$ 把左手电子变为左手正电子。这两个运算中的每一个都把耦合到 $SU(2)$ 规范玻色子的粒子变为不耦合的粒子。然而，这两个运算的组合把左手粒子与右手反粒子互换。因此，组合运算 $CP$ 是 \eqref{eq:20.126} 的对称性。该拉格朗日量在时间反演下也不变。于是我们看到，离散对称性 $C$ 和 $P$ 一方面、$CP$ 和 $T$ 另一方面，在规范场理论中处于非常不同的地位。任何手征规范理论都会自然地破坏 $C$ 和 $P$。在我们分析的这一层面，为什么观测到 $C$ 和 $P$ 是自然界的近似对称性是一个谜。另一方面，每一个规范玻色子与无质量费米子的理论都尊重 $CP$ 和 $T$。实验上已知自然界含有某种破坏 $CP$ 的相互作用，因为 $K^0$ 介子的衰变中 $CP$ 选择定则被微弱破坏。但要为这一破坏寻找来源，我们必须向我们的基本规范理论 \eqref{eq:20.126} 添加项。

首先，我们必须向 \eqref{eq:20.126} 添加会导致 $SU(2)\times U(1)$ 自发破缺的动力学。我们将从含有一个希格斯标量场 $\phi$ 的最简单模型出发。$\phi$ 的最一般的可重整化拉格朗日量为
\begin{equation}
  \mathcal{L}_{\phi} = |D_\mu\phi|^2 + \mu^2\phi^{\dagger}\phi
    - \lambda(\phi^{\dagger}\phi)^2 .
  \label{eq:20.130}
\end{equation}
$\mathcal{L}_{\phi}$ 的厄米性意味着参数 $\mu^2$ 和 $\lambda$ 是实的。因此该拉格朗日量尊重 $P$、$C$ 和 $T$。如上一节末尾所讨论，该拉格朗日量还自动具有产生质量关系 \eqref{eq:20.117} 所需的守衡 $SU(2)$ 对称性。

最后，我们必须添加把希格斯场耦合到夸克和轻子的项。这里，可重整化性和规范不变性提供的约束最弱，因此允许的相互作用有很多。我们将首先分析 $\phi$ 对夸克场的耦合，然后把这一讨论推广到轻子。

在写出希格斯场对夸克的耦合时，我们应该回忆已知存在三代夸克和轻子。因此有三个左手夸克双重态：
\begin{equation}
  Q^i_L = \begin{pmatrix}u_L\\ d_L\end{pmatrix},\qquad
  \begin{pmatrix}c_L\\ s_L\end{pmatrix},\qquad
  \begin{pmatrix}t_L\\ b_L\end{pmatrix} .
  \label{eq:20.131}
\end{equation}
有六个右手夸克，三个具有 $Y = \frac{2}{3}$，三个具有 $Y = -\frac{1}{3}$：
\begin{equation}
  u^i_R = (u_R, c_R, t_R),\qquad
  d^i_R = (d_R, s_R, b_R) .
  \label{eq:20.132}
\end{equation}
当我们把规范场耦合到这些夸克时，我们把普通导数换成协变导数。这自动赋予所有夸克对 QCD 相同的耦合，并赋予所有同类型夸克对弱相互作用和电磁相互作用相同的耦合。

希格斯场对夸克的最一般的可重整化汤川耦合为
\begin{equation}
  \Delta\mathcal{L} = -\lambda^d_{ij}\ \bar{Q}^i_L\cdot\phi\,d^j_R
    - \lambda^u_{ij}\ \epsilon_{ab}\ \bar{Q}^{ia}_L\phi^{b*}\,u^j_R
    + \mathrm{h.c.} ,
  \label{eq:20.133}
\end{equation}
其中 $\lambda^d_{ij}$ 和 $\lambda^u_{ij}$ 是世代空间中的任意复矩阵。在自发对称性破缺之后，这些项产生夸克质量矩阵
\begin{equation}
  \bigl( M_d \bigr)_{ij} = \frac{v}{\sqrt{2}}\lambda^d_{ij},\qquad
  \bigl( M_u \bigr)_{ij} = \frac{v}{\sqrt{2}}\lambda^u_{ij} .
\end{equation}
这些矩阵一般不是对角的。我们可以通过对左手与右手夸克场做幺正变换来把它们对角化：
\begin{equation}
  u_L = U^u_L u'_L,\qquad u_R = U^u_R u'_R,\qquad
  d_L = U^d_L d'_L,\qquad d_R = U^d_R d'_R .
  \label{eq:20.134}
\end{equation}
在这些变换之后，夸克质量矩阵是对角的，具有正特征值。然而，通过 $W$ 玻色子把 $u_L$ 耦合到 $d_L$ 的带电流弱相互作用，在质量基中不再是对角的。混合由 Cabibbo-Kobayashi-Maskawa（CKM）矩阵描述，
\begin{equation}
  V = U^{u\dagger}_L U^d_L .
  \label{eq:20.135}
\end{equation}
CKM 矩阵是幺正 $3\times3$ 矩阵，可以用三个角和一个破坏 $CP$ 的相位参数化。该相位的存在意味着，如果矩阵不是实的，CKM 矩阵就为弱相互作用中的 $CP$ 破坏提供来源。这一相位的存在是标准模型的关键预言之一，也是所观测到的 $K^0$ 系统中 $CP$ 破坏的解释。

如果存在右手中微子，同样的分析也适用于轻子。如果中微子无质量，就没有轻子混合，CKM 矩阵的轻子类似物是平凡的。如果中微子通过 Dirac 或 Majorana 质量项获得质量，就存在非平凡的轻子混合矩阵——Pontecorvo-Maki-Nakagawa-Sakata（PMNS）矩阵。

\problems

\textbf{20.1} \emph{由伴随标量引起的 $SU(3)$ 破缺。} 考虑一个具有伴随表示中标量场的 $SU(3)$ 规范理论，并分析真空期望值 \eqref{eq:20.41} 的对称性破缺模式。显式计算规范玻色子质量，并验证结果 \eqref{eq:20.40} 和 \eqref{eq:20.42}。

\textbf{20.2} \emph{阿贝尔情形的希格斯机制。} 考虑第~20.1 节的阿贝尔 Higgs 模型。给出该理论在幺正规范下的费曼规则，并显式验证戈德斯通玻色子已从谱中被消去，而规范玻色子获得了质量和纵向极化。

\textbf{20.3} \emph{$e^{+}e^{-}$ 前后不对称性。} 计算 GWS 理论中 $e^{+}e^{-}\to f\bar{f}$ 的微分截面，包括 $Z^0$ 交换图的效应。证明在 $Z^0$ 共振峰上，截面具有前后不对称性，其大小由 \eqref{eq:20.80} 中的耦合决定。用参数 $\sin^2\theta_W$ 表示该不对称性。

\textbf{20.4} \emph{中性流中微子散射。} 利用 \eqref{eq:20.80} 中的 $Z^0$ 耦合，计算深非弹中微子散射的中性流贡献。证明中性流与带电流截面之比是 $\sin^2\theta_W$ 的量度，并与实验测定值比较。

\textbf{20.5} \emph{戈德斯通玻色子等价定理。} 在本习题中我们推导戈德斯通玻色子等价定理，它断言在高能下发射纵向 $W$ 或 $Z$ 玻色子的振幅等于发射相应戈德斯通玻色子的振幅。证明在远大于 $m_W$ 的能量下，有质量规范玻色子的纵向极化矢量近似为
\begin{equation}
  \epsilon^\mu_L(k) \approx \frac{k^\mu}{m_W} ,
\end{equation}
并且该纵向模式的耦合等价于相应戈德斯通玻色子的耦合。
